% ======================================================================
% Diagrama de flujo — generado automáticamente por generar_diagrama_flujo.py
% NO EDITAR MANUALMENTE; regenerar desde el descriptor JSON.
% ======================================================================
\documentclass[11pt,a4paper]{article}

% ── Codificación, fuentes y lenguaje ───────────────────────────────────────────
\usepackage[utf8]{inputenc}
\usepackage[T1]{fontenc}
\usepackage[default,scale=0.95]{sourcesanspro}
\renewcommand{\familydefault}{\sfdefault}
\usepackage[spanish,es-noshorthands,es-tabla]{babel}

% ── Geometría y espaciado ─────────────────────────────────────────────────────
\usepackage[top=2.2cm, bottom=2.5cm, left=2.2cm, right=2.2cm, headheight=15pt]{geometry}
\usepackage{setspace}
\setstretch{1.18}
\usepackage{parskip}

% ── TikZ (simbolos ISO 5807 / ANSI X3.5) ──────────────────────────────────────
\usepackage{tikz}
\usetikzlibrary{babel, arrows.meta, positioning, shapes.geometric,
                shapes.symbols, decorations.pathmorphing, calc}

% ── Tablas y listas ───────────────────────────────────────────────────────────
\usepackage{booktabs}
\usepackage{tabularx}
\usepackage{array}
\usepackage{enumitem}

% ── Colores, cajas e iconos ───────────────────────────────────────────────────
\usepackage[dvipsnames]{xcolor}
\usepackage{tcolorbox}
\tcbuselibrary{skins, breakable}
\usepackage{fontawesome5}

% ── Secciones con barra de color (infografía) ─────────────────────────────────
\usepackage{titlesec}
\titleformat{\section}
  {\normalfont\LARGE\bfseries\color{azulNoche}}
  {\colorbox{cyanNeon}{\color{fondoOscuro}\thesection}}{0.7em}{}
  [\vspace{2pt}\color{cyanNeon}\rule{\linewidth}{1.8pt}]
\titlespacing*{\section}{0pt}{24pt}{10pt}
\titleformat{\subsection}
  {\normalfont\large\bfseries\color{azulNoche}}
  {\textcolor{cyanNeon}{\thesubsection}}{0.7em}{}
  [\vspace{1pt}\color{grisLinea}\rule{\linewidth}{0.6pt}]
\titlespacing*{\subsection}{0pt}{14pt}{6pt}

% ── Cabeceras y pies de página ────────────────────────────────────────────────
\usepackage{fancyhdr}
\usepackage{lastpage}
\pagestyle{fancy}
\fancyhf{}
\renewcommand{\headrulewidth}{0pt}
\renewcommand{\footrulewidth}{0pt}
\fancyfoot[C]{%
  \begin{minipage}{\linewidth}
    \vspace{2pt}
    \color{cyanNeon}\rule{\linewidth}{1.2pt}\\[2pt]
    \centering\color{grisTexto}\footnotesize
    Página \thepage\ de \pageref{LastPage}
  \end{minipage}%
}

% ── Hiperenlaces ──────────────────────────────────────────────────────────────
\usepackage[colorlinks=true, linkcolor=azulNoche, urlcolor=cyanNeon]{hyperref}
\sloppy
\emergencystretch 3em

% ── Paleta de colores ─────────────────────────────────────────────────────────
\definecolor{fondoOscuro}{HTML}{0D1B2A}
\definecolor{verdeTurquesa}{HTML}{004D40}
\definecolor{azulNoche}{HTML}{1B3A6B}
\definecolor{azulMedio}{HTML}{2A5298}
\definecolor{cyanNeon}{HTML}{00C8E0}
\definecolor{verdeSignal}{HTML}{00B887}
\definecolor{naranjaVivo}{HTML}{FF7043}
\definecolor{rojoAlerta}{HTML}{E53935}
\definecolor{grisPapel}{HTML}{F4F6FA}
\definecolor{grisLinea}{HTML}{C8D0E0}
\definecolor{grisTexto}{HTML}{6B7A99}
\definecolor{amarilloNota}{HTML}{FFB300}

% ── Tarjetas ──────────────────────────────────────────────────────────────────
\tcbset{
  cajaContenido/.style={
    enhanced, breakable,
    arc=5pt, outer arc=5pt,
    colback=grisPapel, colframe=grisLinea,
    boxrule=0.8pt,
    fonttitle=\bfseries\small, coltitle=azulNoche,
    top=6pt, bottom=6pt, left=10pt, right=10pt,
  },
}

% ── Macros ────────────────────────────────────────────────────────────────────
\newcommand{\iconotexto}[2]{\textcolor{cyanNeon}{\faIcon{#1}}~#2}
\newcommand{\iconoBanda}{sitemap}
\newcommand{\bandaTitulo}[3][fondoOscuro]{%
  \noindent
  \begin{tcolorbox}[%
    enhanced, arc=4mm, colback=#1!10!white, colframe=#1, boxrule=0pt,
    borderline west={6pt}{0pt}{cyanNeon},
    top=6pt, bottom=6pt, left=8pt, right=8pt,
  ]
    \noindent
    \begin{minipage}[c]{0.10\linewidth}
      \centering\color{cyanNeon}\faIcon{clipboard-check}
    \end{minipage}%
    \hspace{8pt}%
    \begin{minipage}[c]{0.88\linewidth}
      {\fontsize{20}{24}\selectfont\bfseries\color{white}#2}\par\vspace{5pt}%
      \textcolor{cyanNeon!80!white}{\small\faIcon{angle-right}~#3}%
    \end{minipage}
    \vspace{4pt}
  \end{tcolorbox}%
  \vspace{8pt}%
}


  % ── Estilos de flujo (ISO 5807 / ANSI X3.5) ─────────────────────────────
  \tikzset{
    flujo/.style={
      >=Stealth,
      font=\footnotesize\sffamily,
    },
    terminador/.style={
      draw=azulNoche, fill=azulNoche!8!white, rounded corners=3pt,
      text width=1.6cm, align=center, inner sep=2mm,
      font=\footnotesize\sffamily,
    },
    proceso/.style={
      draw=azulNoche, fill=cyanNeon!10!white,
      text width=1.6cm, align=center, inner sep=2mm,
      font=\footnotesize\sffamily,
    },
    entradasalida/.style={
      draw=azulNoche, fill=verdeSignal!10!white,
      trapezium, trapezium left angle=75, trapezium right angle=105,
      text width=1.8cm, align=center, inner sep=2mm,
      font=\footnotesize\sffamily,
    },
    decision/.style={
      draw=azulNoche, fill=amarilloNota!20!white,
      diamond, aspect=1.6, text width=1.5cm, align=center, inner sep=1pt,
      font=\footnotesize\sffamily,
    },
    almacenamiento/.style={
      draw=azulNoche, fill=grisPapel,
      cylinder, shape border rotate=90, aspect=0.3,
      text width=1.7cm, align=center, inner sep=2mm,
      font=\footnotesize\sffamily,
    },
    documento/.style={
      draw=azulNoche, fill=azulMedio!8!white,
      text width=1.8cm, align=center, inner sep=2mm,
      font=\footnotesize\sffamily,
    },
    nota/.style={
      draw=grisLinea, dashed, fill=grisPapel,
      text width=1.8cm, align=center, inner sep=2mm,
      font=\scriptsize\sffamily,
    },
    etiqueta/.style={
      font=\scriptsize\sffamily, fill=white,
      fill opacity=0.75, text opacity=1, inner sep=1pt,
    },
  }


% ── Cabecera del documento ────────────────────────────────────────────────────
\fancyhead[L]{\small\color{azulNoche}\textbf{ELECTRONICA}}
\fancyhead[R]{\small\color{grisTexto}flujo\_auditar\_repo}
\renewcommand{\headrulewidth}{0pt}

% ─────────────────────────────────────────────────────────────────────────────
\begin{document}

\bandaTitulo{Auditoria de higiene del repo}{Trece dimensiones, siete nuevas y seis reutilizadas, y un veredicto que gana el peor estado: nunca un promedio}

\vspace{4pt}
\begin{center}
\begin{tcolorbox}[cajaContenido, title={\faIcon{map-signs}~Leyenda de símbolos---ISO 5807 / ANSI X3.5}]
\small
Óvalo = \textbf{terminador} (inicio/fin) $\cdot$
Rectángulo = \textbf{proceso} (acción determinista) $\cdot$
Rombo = \textbf{decisión} (ramas Sí/No) $\cdot$
Paralelogramo = \textbf{entrada/salida} (lectura/escritura de archivos) $\cdot$
Cilindro = \textbf{almacenamiento online} (estado, snapshots) $\cdot$
Rectángulo con base ondulada = \textbf{documento} (sesiones, logs) $\cdot$
Rectángulo punteado gris = \textbf{nota explicativa} (sin acción).
Flujo: arriba$\rightarrow$abajo, izquierda$\rightarrow$derecha.

\faIcon{tags}~Cada símbolo lleva una \textbf{etiqueta} (T/P/D/E/A/N + número, según su tipo);
su significado se expande en la tabla \emph{Leyenda de etiquetas} bajo cada diagrama.
\end{tcolorbox}
\end{center}

\vspace{6pt}

\section{\iconotexto{shield-alt}{La pasada: medir trece dimensiones sin que ninguna oculte a otra}}

\begin{center}
\begin{tikzpicture}[flujo, >=Stealth, x=1cm, y=1cm]
  % Fondo decorativo
  \fill[azulNoche!3!white, rounded corners=4pt]
    (-2.3,1.8) rectangle (5.8,-19.9);
  \node[terminador] (T1) at (0.0,0.0) {T1};
  \node[proceso] (P1) at (0.0,-2.2) {P1};
  \node[decision] (D1) at (0.0,-4.4) {D1};
  \node[terminador] (T2) at (3.5,-4.4) {T2};
  \node[proceso] (P2) at (0.0,-6.6) {P2};
  \node[proceso] (P3) at (0.0,-8.8) {P3};
  \node[decision] (D2) at (0.0,-11.0) {D2};
  \node[nota] (N1) at (3.5,-11.0) {N1};
  \node[proceso] (P4) at (0.0,-13.2) {P4};
  \node[almacenamiento] (A1) at (0.0,-15.4) {A1};
  \node[terminador] (T3) at (0.0,-17.6) {T3};

  \draw[->] (T1.south) -- (P1.north);
  \draw[->] (P1.south) -- (D1.north);
  \draw[->] (D1.east) -- (T2.west) node[etiqueta, above, pos=0.5]{Si};
  \draw[->] (D1.south) -- (P2.north) node[etiqueta, left, pos=0.5]{No};
  \draw[->] (P2.south) -- (P3.north);
  \draw[->] (P3.south) -- (D2.north);
  \draw[->] (D2.south) -- (P4.north) node[etiqueta, left, pos=0.5]{No};
  \draw[->] (D2.east) -- (N1.west) node[etiqueta, above, pos=0.5]{Si};
  \draw[->] (N1.west) -- (P4.east);
  \draw[->] (P4.south) -- (A1.north);
  \draw[->] (A1.south) -- (T3.north);
\end{tikzpicture}
\end{center}

\vspace{6pt}
\begin{tcolorbox}[cajaContenido, title={\faIcon{table}~Leyenda de etiquetas}]
\small
\begin{tabularx}{\linewidth}{@{}>{\centering\arraybackslash}p{1.4cm}X@{}}
\toprule
\textbf{Etiqueta} & \textbf{Significado} \\ \midrule
A1 & Informe de la pasada: .tmp/auditoria\_repo.json, con cobertura explicita de lo medido y lo no medido. \\
D1 & --dimension nombra algo que no existe en el conjunto conocido? \\
D2 & Alguna dimension dio excepcion, se paso de timeout, o su interprete no reconoce la forma de la salida? \\
N1 & Esa dimension cae a no\_verificado y la pasada CONTINUA. Una comprobacion rota no puede tapar a las otras doce ni abortar la auditoria. \\
P1 & Acotar: reunir las 7 dimensiones nuevas y las 6 reutilizadas. Con --rapido se omiten texto, barrera y test\_texto, y lo omitido se declara en vez de contarse como limpio. \\
P2 & Medir las 7 nuevas: auditar\_repo.py comprueba secretos, logs append-only, directivas muertas, capa 3 ausente, peso del indice, ficheros sin trackear y PDF desfasado. \\
P3 & Medir las 6 reutilizadas: texto, estado de sesion, bitacoras, disco, barrera anti-borrado y regresion del verificador de prosa, cada una en subprocess con su timeout. \\
P4 & Traducir: cada dimension queda con su estado y su evidencia contable (ficheros, MB, referencias, hashes, aserciones). \\
T1 & Inicio: ejecutar flujo\_auditar\_repo.py, opcionalmente con --dimension, --json o --rapido. \\
T2 & Salir con codigo 3 por uso incorrecto. No se mide nada: un nombre mal escrito no puede medirse en silencio. \\
T3 & Pasar al algebra del veredicto (seccion siguiente) con la lista de dimensiones ya medida. \\
\bottomrule
\end{tabularx}
\end{tcolorbox}

\section{\iconotexto{gavel}{El veredicto: gana el peor estado, y lo que no se comprobo se dice}}

\begin{center}
\begin{tikzpicture}[flujo, >=Stealth, x=1cm, y=1cm]
  % Fondo decorativo
  \fill[azulNoche!3!white, rounded corners=4pt]
    (-2.3,1.8) rectangle (5.8,-19.9);
  \node[terminador] (T4) at (0.0,0.0) {T4};
  \node[proceso] (P5) at (0.0,-2.2) {P5};
  \node[nota] (N2) at (3.5,-2.2) {N2};
  \node[decision] (D3) at (0.0,-4.4) {D3};
  \node[terminador] (T7) at (3.5,-4.4) {T7};
  \node[decision] (D4) at (0.0,-8.8) {D4};
  \node[terminador] (T8) at (3.5,-8.8) {T8};
  \node[decision] (D5) at (0.0,-13.2) {D5};
  \node[terminador] (T9) at (3.5,-13.2) {T9};
  \node[terminador] (T6) at (0.0,-15.4) {T6};
  \node[almacenamiento] (A2) at (0.0,-17.6) {A2};
  \node[nota] (N3) at (3.5,-17.6) {N3};

  \draw[->] (T4.south) -- (P5.north);
  \draw[->] (P5.south) -- (D3.north);
  \draw[->] (D3.east) -- (T7.west) node[etiqueta, above, pos=0.5]{Si};
  \draw[->] (D3.south) -- (D4.north) node[etiqueta, left, pos=0.5]{No};
  \draw[->] (D4.east) -- (T8.west) node[etiqueta, above, pos=0.5]{Si};
  \draw[->] (D4.south) -- (D5.north) node[etiqueta, left, pos=0.5]{No};
  \draw[->] (D5.east) -- (T9.west) node[etiqueta, above, pos=0.5]{Avisos};
  \draw[->] (D5.south) -- (T6.north) node[etiqueta, left, pos=0.5]{Ninguno};
\end{tikzpicture}
\end{center}

\vspace{6pt}
\begin{tcolorbox}[cajaContenido, title={\faIcon{table}~Leyenda de etiquetas}]
\small
\begin{tabularx}{\linewidth}{@{}>{\centering\arraybackslash}p{1.4cm}X@{}}
\toprule
\textbf{Etiqueta} & \textbf{Significado} \\ \midrule
A2 & En cualquiera de las cuatro salidas: flujo/inicio y flujo/fin en .tmp/session\_log\_<run\_id>.jsonl, y el estado derivado en .tmp/run\_state.json CON run\_id (sin el, la vista se clasifica como corrupta). \\
D3 & Hay alguna dimension en estado fallo? \\
D4 & Sin fallos: hay alguna dimension en no\_verificado? \\
D5 & Todas medidas y sin fallos: hay avisos? \\
N2 & Sin media y sin pesos. Doce dimensiones sanas sobre trece darian un promedio de 0.92 y taparian un disco al 91\%: el promedio esconde justo lo que hay que arreglar. \\
N3 & La auditoria no modifica nada: corregir lo que encontro (sacar bloat del indice, purgar vistas huerfanas, recompilar PDF) son decisiones del operador, con sus consecuencias. \\
P5 & Contar estados. Un estado fuera del vocabulario se degrada a no\_verificado y nombra la dimension culpable. \\
T4 & Entrada del algebra: la lista de dimensiones con su estado y su evidencia contable. \\
T6 & Salir con codigo 0 y veredicto limpio, y solo cuando la cobertura fue completa. Un verde sin cobertura es la peor respuesta posible. \\
T7 & Salir con codigo 1, veredicto con\_fallos, e imprimir las acciones en orden de gravedad: primero los fallos, despues los avisos. \\
T8 & Salir con codigo 2 y enumerar que dimension no se pudo comprobar y por que. Sin verificar no es un veredicto: es una laguna declarada. \\
T9 & Salir con codigo 0 y veredicto con\_avisos: cada aviso queda nombrado y con su accion, sin bloquear. \\
\bottomrule
\end{tabularx}
\end{tcolorbox}


\section{\iconotexto{sticky-note}{Notas}}
\begin{itemize}
\item El veredicto global es el PEOR estado de las dimensiones. No hay puntuacion ni media: 12 dimensiones sanas sobre 13 darian 0.92 y taparian un disco al 91\%
\item No verificado no es estar bien. Una dimension que no se pudo medir impide que el conjunto salga como limpio y devuelve codigo 2
\item Un estado fuera del vocabulario \{ok, aviso, fallo, no\_verificado\} se degrada a no\_verificado: una errata en una cadena no puede producir un verde
\item La auditoria es de solo lectura. Escribe unicamente sus tres salidas en .tmp/ y no borra, no toca el indice de git ni reescribe historia
\item Ningun estado de la sesion se borra ni se toca: la limpieza de vistas huerfanas la decide el operador con estado\_sesion.py clean
\item Cero gasto de LLM y cero red: las trece comprobaciones son codigo determinista y no pasan por el enrutador
\end{itemize}

% ─────────────────────────────────────────────────────────────────────────────
\end{document}
